\documentclass[sigchi-a, authorversion]{acmart}
\usepackage{booktabs} % For formal tables
\usepackage{ccicons}  % For Creative Commons citation icons

% Copyright
%\setcopyright{none}
\setcopyright{acmcopyright}
%\setcopyright{acmlicensed}
%\setcopyright{rightsretained}
%\setcopyright{usgov}
%\setcopyright{usgovmixed}
%\setcopyright{cagov}
%\setcopyright{cagovmixed}


% DOI
%\acmDOI{10.475/123_4}

% ISBN
%\acmISBN{123-4567-24-567/08/06}

%Conference
\acmConference[SPWAL LMU]{Wissenschaftliches Arbeiten und Lehren, LMU}{Novemver 2017}{Munich, Germany} 
%\acmYear{1997}
%\copyrightyear{2016}

%\acmPrice{15.00}

%\acmBadgeL[http://ctuning.org/ae/ppopp2016.html]{ae-logo}
%\acmBadgeR[http://ctuning.org/ae/ppopp2016.html]{ae-logo}

\begin{document}
\title{Identifying the challenges towards distributed nano data center infrastructor}

\author{Katrin Kolb}
\affiliation{%
  \institution{Ludwig Maximilian University of Munich}
  \city{Munich} 
  \country{Germany}}
\email{katrinkolb@web.de}

\author{Katharina Rupp}
\affiliation{%
  \institution{Ludwig Maximilian University of Munich}
  \city{Munich} 
  \country{Germany}}
\email{katharina.rupp@web.de}

\author{Mengchu Li}
\affiliation{%
  \institution{Ludwig Maximilian University of Munich}
  \city{Munich} 
  \country{Germany}}
\email{mengchu.li@yahoo.com}

\author{Melanie Hauser}
\affiliation{ %
  \institution{Ludwig Maximilian University of Munich}
  \city{Munich}
  \country{Germany}}
\email{Melanie.Hauser@campus.lmu.de}

\author{Andreas Scholz}
\affiliation{%
  \institution{Ludwig Maximilian University of Munich}
  \city{Germany}
  \country{Germany}}
\email{Andreas.Scholz@campus.lmu.de}

\author{Diana Irmscher} 
\affiliation{%
 \institution{Ludwig Maximilian University of Munich}
 \country{Germany}}
\email{d.irmscher@campus.lmu.de}

% The default list of authors is too long for headers.
\renewcommand{\shortauthors}{Katrin Kolb et al.}

%
% The code below should be generated by the tool at
% http://dl.acm.org/ccs.cfm
% Please copy and paste the code instead of the example below. 
%
\begin{CCSXML}
<ccs2012>
 <concept>
 <concept_id>10003033.10003106.10003110</concept_id>
 <concept_desc>Networks~Data center networks</concept_desc>
 <concept_significance>500</concept_significance>
 </concept>
 <concept>
 <concept_id>10003456.10003457.10003490.10003507.10003508</concept_id>
 <concept_desc>Social and professional topics~Centralization / decentralization</concept_desc>
 <concept_significance>500</concept_significance>
 </concept>
 <concept>
 <concept_id>10011007.10010940.10010971.10011120</concept_id>
 <concept_desc>Software and its engineering~Distributed systems organizing principles</concept_desc>
 <concept_significance>300</concept_significance>
 </concept>
 <concept>
 <concept_id>10010520.10010521.10010537.10003100</concept_id>
 <concept_desc>Computer systems organization~Cloud computing</concept_desc>
 <concept_significance>100</concept_significance>
 </concept>
 <concept>
 <concept_id>10010520.10010521.10010537.10010540</concept_id>
 <concept_desc>Computer systems organization~Peer-to-peer architectures</concept_desc>
 <concept_significance>100</concept_significance>
 </concept>
% <concept>
%  <concept_id>10010520.10010575.10010755</concept_id>
%  <concept_desc>Computer systems organization~Redundancy</concept_desc>
%  <concept_significance>300</concept_significance>
% </concept>
% <concept>
%  <concept_id>10010520.10010553.10010554</concept_id>
%  <concept_desc>Computer systems organization~Robotics</concept_desc>
%  <concept_significance>100</concept_significance>
% </concept>
% <concept>
%  <concept_id>10003033.10003083.10003095</concept_id>
%  <concept_desc>Networks~Network reliability</concept_desc>
%  <concept_significance>100</concept_significance>
% </concept>
</ccs2012>
\end{CCSXML}
\ccsdesc[500]{Networks~Data center networks}
\ccsdesc[500]{Social and professional topics~Centralization / decentralization}
\ccsdesc[300]{Software and its engineering~Distributed systems organizing principles}
\ccsdesc[100]{Computer systems organization~Cloud computing}
\ccsdesc[100]{Computer systems organization~Peer-to-peer architectures}

%\ccsdesc[300]{Computer systems organization~Redundancy}
%\ccsdesc{Computer systems organization~Robotics}
%\ccsdesc[100]{Networks~Network reliability}


\begin{abstract} % Andreas
  UPDATED---\today. 
\end{abstract}


\keywords{Authors' choice; of terms; separated; by
  semicolons; include commas, within terms only; required.}

\maketitle

%\begin{sidebar}
%  \textbf{Good Utilization of the Side Bar} 
%  
%  \textbf{Preparation:} Do not change the margin
%  dimensions and do not flow the margin text to the
%  next page. 
%  
%  \textbf{Materials:} The margin box must not intrude
%  or overflow into the header or the footer, or the gutter space
%  between the margin paragraph and the main left column. 
%  
%  \textbf{Images \& Figures:} Practically anything
%  can be put in the margin if it fits. Use the
%  \texttt{{\textbackslash}marginparwidth} constant to set the
%  width of the figure, table, minipage, or whatever you are trying
%  to fit in this skinny space.
%
%  \caption{This is the optional caption}
%  \label{bar:sidebar}
%\end{sidebar}

%\begin{figure}
%  \includegraphics[width=\marginparwidth]{sigchi-logo}
%  \caption{Insert a caption below each figure.}
%  \label{fig:sample}
%\end{figure}


\section{Introduction} % Katrin

\section{Related Work} % Melanie, Katharina
Some concepts in the field of nano data centers have been developed. For example, Valancious et al. introduced NaDa. NaDa is a distributed computing platform, which uses a managed peer-to-peer model for its infrastructure. They furthermore evaluated their system in terms of energy savings. \\
Jalali et al. \\
Laoutaris et al. \\

\section{Justification} % Melanie, Katharina
Although a lot of research has been done and concepts have been developed, nano data centers have not yet actually been implemented in big extent. To realize nano data centers, it is important to know why it has not been done yet and which obstacles have to be resolved. That is why this paper analyses the causes and obstacles.

\section{Evaluation} % Mengchu

\paragraph{Expected achievements of the project:} A list of challenges towards the development of nano data center.
We will first analyse the features that are related to nano data center development and then list the challenges that need to be overcome. 
For a feature to be listed as a challenge, 
the following conditions must be satisfied:
\begin{enumerate}
\item[1.] The feature is a necessary prerequisite for the development of nano data center;
\item[2.] the current status of the feature do not meet the demand of nano data center development. 
\end{enumerate}

\paragraph{Evaluation method:} 
\begin{enumerate}
\item[1.] To test the first condition, 
we will study the existing nano data center models proposed by other research, 
and find out how our proposed challenges are involved in these models. 
For example, whether a proposed challenge is related to the components that construct the infrastructure of these nano data centers,
or which functionality supported by these nano data centers will be influenced by the challenge?
\item[2.] To test the second condition,
we will formulate a report of the current status of the proposed challenges,
and compare the results with their expected status derived from the data center models proposed by other research.
If the current status does not match the expected status,
we will try to find out the reason and propose some approaches to narrow the gap.
\end{enumerate}

\paragraph{Software and resources for the evaluation:}
According to the current evaluation plan, no software needs to be built. Paper survey will be essential for carrying out the evaluation.

\section{Research Plan} % Diana
\begin{table}[H]
  \caption{Research Plan}
  \label{tab:researchPlan}
  \begin{tabular}{ll}
    \toprule
    until 06.11.2017 & Preparation and submission of research proposal \\
    & Literature review to determine SRP's context \\
    & Narrow the scope of the project \\
    \hline
    07.11.2017 & Presentation of the research proposal \\
    & Meeting with the instructor \\
    \hline
    08.11.2017 & Starting project work: \\
    & Research on current nano data centers in practical use \\
    & What has been treated theoretically and practically as nano data centers \\
    & (projects, etc.) \\
    \hline
    until 20.11.2017 & Preparation and submission of progress report I \\
    \hline
    21.11.2017 & Presentation of the progress report I \\
    & Meeting with the instructor \\
    \hline
    22.11.2017 & Continue the project work: \\
    & Research on current nano data centers in practical use \\
    & What has been treated theoretically and practically as nano data centers \\
    & (projects, etc.) \\
    \hline
    12.12.2017 & Mid term synchronisation \\
    & Meeting with the instructor \\
    \hline
    13.12.2017 & Continue the project work: \\
    & Search for obstacles, why the examined works have not yet \\
    & been practically implemented or applied\\
    & Evaluation of the search results \\
    \hline
    until 12.01.2018 & Preparation and submission of progress report II \\
    \hline
    13.01.2018 & Continue the project work: \\
    & Evaluation of the search results \\
    \hline
    until 26.01.18 & Preparation and submission of final deliverables \\
    \hline
    30.01.18 or & Presentation of final deliverables \\
    06.02.18 &  \\
    \bottomrule
  \end{tabular}
\end{table}

\section{Risk Analysis} % Diana
Das ist ein Text \cite{DBLP:conf/conext/ValanciusLMDR09} und das auch \cite{DBLP:journals/sigmetrics/JalaliAVHAT14}

% References
\bibliography{sigchi-a}
\bibliographystyle{ACM-Reference-Format}

\end{document}
